% Job posting: https://jobs.ashbyhq.com/cohere/2d4d224f-94d4-40ce-88d5-5d340bc0da4e
% =====================================================================
%  CV — Nishal Stanislaus Silva, PhD
%  Variant: V1 (AI-Industry) — Cohere Senior Research Engineer, Safety Tooling & Data:
%  data pipelines + real/synthetic data + benchmarking + infra ownership at a frontier lab.
%  Changelog: V1 chosen (research-engineer at a company, not academic). Groups Data & ML Tooling /
%  ML Engineering / Research & Leadership. Portfolio 3 items. Pubs = 3. Real+synthetic data
%  framing from McGill. Metrics 14 ms / F1 0.76 / 300%. Page target 2. No forced \newpage.
% =====================================================================
\documentclass[11pt]{article}
\usepackage[left=0.7in,top=0.7in,right=0.7in,bottom=0.7in]{geometry}
\usepackage{enumitem}
\usepackage[hidelinks]{hyperref}
\usepackage{titlesec}
\usepackage{array}
\usepackage{setspace}
\usepackage[T1]{fontenc}
\usepackage{newtxtext,newtxmath}
\ifdefined\pdfgentounicode
  \input{glyphtounicode}
  \pdfgentounicode=1
\fi
\setstretch{1.05}
\pagenumbering{gobble}
\titleformat{\section}{\Large\scshape}{}{4em}{}[\vspace{-0.7em}\rule{\linewidth}{0.4pt}\vspace{-0.4em}]
\titlespacing*{\section}{0pt}{1em}{0.4em}
\newenvironment{cvrole}[5]{%
  \par\noindent\textbf{\large #1} | \textit{\small #2, #3}\hfill \textit{\small #4 -- #5}\par
  \begin{itemize}[leftmargin=2em, itemsep=-0.3em, topsep=-0.5em]%
}{%
  \end{itemize}\par\noindent\mbox{}\par\vspace{0.1em}%
}
\newcommand{\cvName}{Nishal Stanislaus Silva}
\newcommand{\cvStatus}{Canadian Permanent Resident, Eligible to work in Canada without sponsorship}
\newcommand{\cvTagline}{Machine Learning Research Engineer | Data Pipelines, Benchmarking \& Real-Time Inference}

\begin{document}
\pagestyle{empty}
\begin{center}
    {\LARGE \scshape \textbf{\cvName}}{\large \textit{, PhD}}
    \\[1pt]
    {\small \cvTagline}
    \\[1pt]
    {\footnotesize\textit{\cvStatus}}
    \\[1pt]
    {\small
    \href{mailto:nishal.silva@hotmail.com}{nishal.silva@hotmail.com}
    \,\,|\,\, +1 613 200 2030
    \,\,|\,\, Toronto, ON
    \,\,|\,\, \href{https://nishal.xyz}{nishal.xyz}
    \,\,|\,\, \href{https://www.linkedin.com/in/nishal-silva}{linkedin.com/in/nishal-silva}
    \,\,|\,\, \href{https://github.com/ns2max}{github.com/ns2max}}
\end{center}

\section*{Professional Summary}

Machine Learning Research Engineer with a PhD (University of Trento, 2025) and 10+ years across industry and academia, working at the intersection of data tooling, benchmarking and production ML. Owns data pipelines end to end, acquisition, feature engineering, real and synthetic data integration, training/evaluation and deployment, and designs the benchmark suites (baselines + ablations) that tell teams which architecture and data choices actually move the needle. Postdoctoral work on the EU-funded MUSMET project delivered real-time inference pipelines at 14 ms latency and F1 0.76, with reproducible, documented tooling built for long-term maintainability. Earlier, at MAS Holdings, built real-time ETL and ingestion processing millions of IoT events per day across distributed sites, and set the CI/CD and code-review standards that scaled ML adoption across engineering teams. Published researcher (10 peer-reviewed papers) with a benchmark-first, statistically rigorous approach to experimental design.

\section*{Core Competencies}

\textbf{Data \& ML Tooling:} data pipeline design · real + synthetic data integration · ETL pipelines · data validation and quality processes · benchmarking and ablations · experimental design · reproducible tooling\\
\textbf{ML Engineering:} end-to-end ML lifecycle · real-time / low-latency inference · production deployment · MLOps · human-in-the-loop systems · model compression\\
\textbf{Research \& Leadership:} statistical analysis · algorithm design · cross-functional collaboration · technical-standards ownership · mentoring · 10+ peer-reviewed publications

\section*{Technical Stack}

\textbf{Languages:} Python · C++ · C · JavaScript · MATLAB · Shell\\
\textbf{ML / AI:} PyTorch · TensorFlow, Keras · scikit-learn · HuggingFace Transformers · ONNX · SageMaker · transfer learning · diffusion models · VAE-based synthesis\\
\textbf{Data \& Analytics:} SQL · Snowflake · AWS (EC2, S3, ECS, MWAA, RDS) · NumPy, Pandas, SciPy · ETL pipelines · information retrieval\\
\textbf{Dev \& MLOps:} Docker · Kubernetes · MLflow · Airflow · Jenkins · Git · pytest · CI/CD · REST APIs

\section*{Education}
\textbf{Ph.D -- Information Engineering and Computer Science} \textit{\small{ -- University of Trento, Italy}} \hfill \textit{Jan 2025}\\[2pt]
\noindent\textbf{M.Sc -- Telecommunication and Electronic Engineering} \textit{\small{ -- Sheffield Hallam University, UK}} \hfill \textit{Aug 2018}\\[2pt]
\noindent\textbf{B.Eng (Hons) -- Electronic Engineering} \textit{\small{ -- Sheffield Hallam University, UK}} \hfill \textit{Mar 2014}

\section*{Portfolio and Demos \footnotesize {\normalfont{(full portfolio: \href{https://nishal.xyz/research}{\textit{nishal.xyz/research}})}}}
\begin{itemize}[leftmargin=1.2em, itemsep=-0.25em]
    \item \href{https://github.com/ns2max/nebula}{\textbf{Nebula}} -- open-source C++17 feature-extraction library with per-feature ARM latency benchmarks.
    \item \href{https://zenodo.org/record/10818617}{\textbf{DoMP / DoPP / DoDP / DoDP2}} -- four open datasets, 9,800+ recordings, designed as worst-case benchmarks for real-time pattern detection.
    \item \href{https://marketplace.visualstudio.com/items?itemName=ns2max.livelatex}{\textbf{LiveLaTeX}} -- published VS Code extension (real-time compile and preview).
\end{itemize}

\section*{Professional Experience}

\begin{cvrole}{Postdoctoral Researcher}{University of Trento}{Trento, Italy}{Jan 2025}{Dec 2025}
    \item Owned the full data pipeline for streaming audio, sensor and gesture data, acquisition through edge deployment and monitoring, building the tooling that let the team iterate quickly with reproducible results; EU-funded MUSMET (EIC Pathfinder).
    \item Designed benchmark suites (baselines + ablations) to identify architecture and feature-coverage gaps and quantify trade-offs under latency and compute constraints; delivered 14 ms inference at F1 0.76, 74$\times$ faster than the DTW baseline (JAES 2026).
    \item Translated research prototypes into maintainable, production-grade components with reproducible pipelines and deployment docs; standards adopted by academic and industry partners.
\end{cvrole}

\begin{cvrole}{Visiting Researcher}{McGill University}{Montreal, Canada}{May 2024}{Aug 2024}
    \item Built synthetic data generation pipelines (diffusion- and VAE-based) combined with real multimodal sensor and time-series data to improve model robustness under limited-label regimes.
    \item Power and compute profiling to assess real-time inference feasibility; findings shaped deployment decisions.
\end{cvrole}

\begin{cvrole}{Technical Consultant}{Forestpin (Pvt) Ltd}{Colombo, Sri Lanka}{Aug 2020}{Feb 2021}
    \item Designed and deployed ML and statistical pipelines for financial forensics, compliance monitoring and risk detection on large structured enterprise datasets.
    \item Built Python + SQL backend scoring services for anomaly flagging and production triage; translated compliance requirements into technical specs.
\end{cvrole}

\begin{cvrole}{Research Engineer}{MAS Holdings (Pvt) Ltd}{Colombo, Sri Lanka}{Jan 2015}{Jul 2020}
    \item Built real-time ingestion and ETL pipelines processing millions of IoT and operational time-series events per day across distributed sites, with validation and monitoring for consistent reliability.
    \item Led end-to-end ML and CV system development integrating camera and sensor streams into production workflows; operational efficiency improved up to 300\%, inspection time cut 99.5\%.
    \item Established CI/CD and code-review standards, mentored engineers, and scaled ML adoption across teams, turning opinionated engineering practice into team-wide defaults.
\end{cvrole}

\section*{Selected Publications \footnotesize {\normalfont{(full list: \href{https://nishal.xyz/publications}{\textit{nishal.xyz/publications}})}}}
\begin{itemize}[leftmargin=1.2em, itemsep=-0.25em]
    \item Silva, N. and Turchet, L. \textit{Real-Time Audio Pattern Detection for Smart Musical Instruments}. Journal of the Audio Engineering Society, Vol.~74, No.~3, Mar. 2026. \textit{(F1 0.76 at 14 ms on RPi4, 74$\times$ faster than DTW)}
    \item Silva, N. and Turchet, L. \textit{A Structural Similarity Index Based Method to Detect Symbolic Monophonic Patterns in Real-Time}. 25th Int. Conference on Digital Audio Effects (DAFx), Sep. 2022. \textit{(95\% detection, training-free)}
    \item Silva, N., Boem, A. and Turchet, L. \textit{Interactive IoMusT-Based Concerts: Real-Time Pattern Recognition and Audience Experience}. IEEE International Conference on Immersive and 3D Audio (I3DA), Sep. 2025.
\end{itemize}

\vspace{-1em}
\section*{Independent Research, Highlights, \& Recognition}
\begin{itemize}[leftmargin=1.2em, itemsep=-0.25em]
    \item \textbf{Open datasets:} four open-access datasets (DoMP, DoPP, DoDP, DoDP2), 9,800+ recordings from 70+ musicians, published on Zenodo with a shared benchmark design.
    \item \textbf{MIDI Innovation Awards:} Finalist (Sep 2023) -- real-time symbolic pattern detection for smart instruments.
    \item \textbf{GMOA Recognition (Sri Lanka):} computer-vision system for COVID-19 patient vitals monitoring deployed across hospital wards. Reviewer, IEEE Access.
\end{itemize}

\end{document}
